\section{Finite Newton data on the real observation graph}
\label{sec:finite-newton-v95}
Let $X$ be compact semialgebraic, let $F:X\to\mathcal P$ be continuous
semialgebraic, and let $A_x$ be a continuous semialgebraic monic polynomial
of degree $n$. Fix $P\in F(X)$ and suppose $A_x=A_0$ on $F^{-1}(P)$.
The competitors in this section range over all of $X$. A restriction to a
compact local model will always be denoted $X_{\rm loc}$ instead.
Write $A_0(z)=\prod_r(z-r)^{m_r}$. Near the exact fibre, unique monic
cluster factors give
\[
 A_{r,x}(r+z)=z^{m_r}+\sum_{a<m_r}c_{r,a}(x)z^a,\qquad k_{r,a}=m_r-a.
\]
The factor equations, with isolating discs fixed, define semialgebraic
functions by unique choice. Their analyticity in the original polynomial
coefficients follows from the invertible multiplication differential.
These are separate assertions; analyticity alone would not imply
semialgebraicity.

Put $R_{je}=\sqrt{w_j}(\sqrt{F(x)_{je}}-\sqrt{P_{je}})$ and introduce the
compact graph of $x,F,A,(c_{r,a}),R$ over a sufficiently small closed
observation neighbourhood. On this graph set
\begin{equation}\label{eq:ideals-v95}
 G=\sum_{j,e}R_{je}^2=h_w(F(x),P)^2,
 \qquad H_{r,a}=(\Re c_{r,a})^2+(\Im c_{r,a})^2.
\end{equation}
Thus $G$ and $H_{r,a}$ are polynomial functions on the graph coordinates,
even at zero observed cells. We use their real graph, not its unrestricted
complexification.

\begin{definition}[Real monomial data]\label{def:real-data-v95}
A real monomial cover of the graph near $G=0$ is a finite collection of
relatively compact coordinate boxes on smooth real manifolds, with maps
into the graph, which cover an observation neighbourhood and in which
\begin{equation}\label{eq:monomial-data-v95}
 G=u\prod_i|y_i|^{a_i},\qquad
 H_{r,a}=v_{r,a}\prod_i|y_i|^{b_{r,a,i}}.
\end{equation}
A displayed $H$ may instead be identically zero on a box. Nonzero units
are positive and bounded above and below on the closures of smaller
boxes. The exponents are nonnegative integers. Only coordinate divisors
with $a_i>0$ and a real generic point in a box are recorded. In
particular every recorded transversal is an admissible real arc. The
maps need not be injective, and all lower-dimensional pieces are included.
\end{definition}

\begin{lemma}[Construction of the data]\label{lem:real-cover-v95}
The graph in \eqref{eq:ideals-v95} admits real monomial data. With real
algebraic input it admits an effective finite construction. The input
consists of a polynomial formula for $X$ and the graphs of $F,A$, the
algebraic datum $P$, the weights, and an isolating description of the
base clusters. No observation-ball maximum is an input to the construction.
\end{lemma}
\begin{proof}
We spell out the real-locus issue. Uniformize the compact real graph by a
proper map from smooth real manifolds. Simultaneously make the nonzero
functions $G,H_{r,a}$ normal crossings on each component. A finite cover
of the inverse image of $G=0$, followed by shrinking boxes, gives
\eqref{eq:monomial-data-v95}. Components on which a function vanishes
identically are treated separately, not included in a product of
nonzero functions. Properness and compactness ensure that the images of
the boxes contain every sufficiently small observation ball. All source
points map into the original graph, including its inequalities. This is
the uniformization and normal-crossing construction of
\cite{v95BM88}, not an additional resolution theorem.

For effective algebraic input one can avoid an arbitrary analytic
uniformization. Decompose the compact semialgebraic graph into finitely
many basic closed sets. Replace each inequality $p\ge0$ by
$p=s^2$. The slack variables are bounded over the compact graph, so this
produces compact real algebraic lifts. Resolve the lifts over the real
algebraic numbers, and resolve their singular algebraic subsets
recursively when real points are not covered by the resolved regular
locus. The uncovered locus has strictly smaller dimension at each such
step. The real images remain in the lift throughout. This gives a finite
proper cover of its real points. Simultaneous principalization on each
smooth source, applied to the product of the nonzero $G,H_{r,a}$, gives
the required monomials. Constructive characteristic-zero resolution
\cite{v95BM97} supplies the algebraic operations. Real feasibility of a
chart and its divisors is decided by quantifier elimination. One must
retain this test: a divisor seen only on the complexification supplies
no admissible transversal.
\end{proof}

\begin{theorem}[Finite real-valuative Newton formula]
\label{thm:finite-newton-v95}
For each cluster coefficient take the minimum over all real monomial
boxes and their recorded divisors:
\begin{equation}\label{eq:finite-edge-v95}
 \beta_{r,a}=\min\frac{b_{r,a,i}}{a_i},\qquad
 \alpha=\min_{r,a}\frac{\beta_{r,a}}{k_{r,a}}.
\end{equation}
Identically zero functions contribute $+\infty$. Then either all
coefficients vanish on an observation neighbourhood, or
$\alpha\in\mathbb Q_{>0}$ and
\begin{equation}\label{eq:finite-law-v95}
 \omega_{P,X}(t)=C_{P,X}t^\alpha+o(t^\alpha),\qquad 0<C_{P,X}<\infty.
\end{equation}
The numbers $\beta_{r,a}$ are the orders of the corresponding coefficient
envelopes, but formula \eqref{eq:finite-edge-v95} computes them from
finite integer data without constructing those envelopes. Each finite
minimum is attained by an admissible transversal on the resolved real
graph. The formula is independent of the chosen cover. It applies to
singular observation maps, positive-dimensional exact parameter fibres,
and inequality faces of the entire closed stochastic model.
\end{theorem}
\begin{proof}
Fix a nonzero $H=H_{r,a}$ and let $\beta$ be the first minimum in
\eqref{eq:finite-edge-v95}. On every box, $b_i\ge\beta a_i$ for all
$a_i>0$; for $a_i=0$ one has $b_i\ge0$. Bounded units and
$|y_i|\le1$ therefore give $H\le C G^\beta$. A finite cover makes the
constant uniform over the observation neighbourhood. Hence
$|c_{r,a}|\le C h_w^\beta$.

Choose a divisor realizing the minimum and a real point on it away from
all other divisors. Fix the other coordinates at this point and let
$y_i=s>0$. Its image is an admissible arc with
$G\asymp s^{a_i}$ and $H\asymp s^{b_i}$. Inverting the first relation
(or restricting its analytic branch) gives a lower bound of order
$t^\beta$ for the coefficient envelope. This proves the first equality
of orders without optimizing over arcs or balls. Since $H=0$ on $G=0$,
a nonzero $H$ has $b_i>0$ on every recorded observation divisor. Thus
all finite orders are positive. If no nonzero $H$ occurs over such a
divisor, properness gives an observation neighbourhood on which $H=0$.

For a monic cluster polynomial with coefficients $c_a$, its largest root
modulus is comparable, with constants depending only on the degree, to
$\max_a|c_a|^{1/(m-a)}$. This follows from elementary symmetric functions
and the geometric-series root bound. The base spectrum is in every
ball; therefore the diameter is between its maximal displacement and
twice that displacement. This gives the exponent in
\eqref{eq:finite-edge-v95}. The joint coefficient specialization below,
as proved in Proposition~\ref{prop:effective-fibre-v95}, supplies the leading constant and
the little-oh term. Since the envelope orders are intrinsic to the real
observation--spectral graph, two covers give the same minimum.
\end{proof}

\paragraph{What the valuation formula does and does not assert.}
Ratios of vanishing orders on a principalization are classical
\L{}ojasiewicz data; compare \cite[Proposition 4.3 and Theorem 4.8]{v95BE09}
and the broader filtration framework of \cite{v95Ha26}. The additional
steps here are the real probability-graph lift, the cluster weights
$k_{r,a}$, and the joint leading spectral set. Neither bare channel ranks
nor root multiplicities replace the exponents $a_i,b_{r,a,i}$. The
finite cutoff is the output of a resolution of the given graph, not a
universal bound on ordinary Taylor jets.

\subsection{The joint leading set as a real weighted specialization}
Write $\alpha=p/q$ in lowest positive terms. Use $\varepsilon>0$ and
$t=\varepsilon^q$. In the same graph impose
\begin{equation}\label{eq:weighted-graph-v95}
 R_{je}=\varepsilon^q U_{je},\qquad
 c_{r,a}=\varepsilon^{p k_{r,a}}V_{r,a},\qquad
 \sum_{j,e}U_{je}^2\le1.
\end{equation}
Let $\mathscr I_0$ be the $(U,V)$-projection of the fibre
$\varepsilon=0$ of the real closure of this set. The model coordinates
are retained until the closure is taken. They are bounded, so projecting
them afterwards loses no limiting point. Set
$\mathscr C_0=\operatorname{pr}_V\mathscr I_0$ and
\[
 Z_r(V)=Z\left(z^{m_r}+\sum_{a<m_r}V_{r,a}z^a\right).
\]
The simultaneous observation variables in \eqref{eq:weighted-graph-v95}
retain the metric; the simultaneous coefficient variables retain all
relations within and between clusters.

\begin{proposition}[Effective leading fibre and diameter]
\label{prop:effective-fibre-v95}
The set $\mathscr C_0$ is compact and equals the joint limiting
coefficient set of the observation balls. In particular
\begin{equation}\label{eq:leading-diameter-v95}
 C_{P,X}=\max_{V,V'\in\mathscr C_0}\max_r
                    d_\infty(Z_r(V),Z_r(V')).
\end{equation}
For real algebraic input, a finite algorithm outputs a quantifier-free
formula for $\mathscr I_0$, hence for $\mathscr C_0$, and an isolating
polynomial and rational interval for the positive real algebraic number
$C_{P,X}$. It also outputs the rational exponent from
\eqref{eq:finite-edge-v95}. These outputs require no coefficient-envelope
maximization or Puiseux analysis of such a maximum.
\end{proposition}
\begin{proof}
The monomial bound makes $V$ uniformly bounded in
\eqref{eq:weighted-graph-v95}; $U$ already lies in the unit ball.
For every $t>0$ the projected $V$-section is exactly the rescaled
coefficient image of $B_P(t)$. Thus its closure fibre is the joint
coefficient limit. Definability rules out oscillatory subsequential
limits: for a point of that fibre, its distance to the section is a
bounded semialgebraic function with a zero subsequential limit, hence
has limit zero. A finite net gives Hausdorff convergence. Continuity of
monic root multisets on compact coefficient sets gives
\eqref{eq:leading-diameter-v95}. Zero belongs to $\mathscr C_0$ because
an exact-fibre representative belongs to every observation ball.

We specify the output procedure. Encode each algebraic constant by its
minimal polynomial and isolating interval. Represent complex quantities
by real and imaginary coordinates. Isolate each cluster using disjoint
algebraic discs and the polynomial multiplication equations; their
unique near-base branch is selected by strict distance inequalities.
All powers in \eqref{eq:weighted-graph-v95} are integral. Closure is
encoded by the first-order condition that every positive-radius ball
meets the set with $\varepsilon>0$. Real quantifier elimination removes
these quantifiers and the bounded model coordinates. Encode the roots
of a monic polynomial by their elementary symmetric identities. For
each of the finitely many permutations, squared Euclidean distances
encode the bottleneck metric using finite maxima and minima. A final
compact maximization formula isolates the unique value in
\eqref{eq:leading-diameter-v95}. A singleton semialgebraic over the real
algebraic numbers is real algebraic, giving the stated number output.
There is no complexity claim. Algebraic saturation by $\varepsilon$
alone is not this algorithm: its special fibre can contain inadmissible
real points. The real closure and the original inequalities are essential.
\end{proof}

\begin{example}\label{ex:singular-edge-v95}
For $x\in[-1/2,1/2]$, take
$F(x)=(1/2+x^4,1/2-x^4)$ and $A_x(z)=z^2-x^6$, at $x=0$.
The observation derivative vanishes. Here $G$ has order eight and the
only nonzero $H$ has order twelve. Thus $\beta=3/2$ and $\alpha=3/4$.
This is a finite singular-chart calculation, not the immersed-jet
formula. It also illustrates why observation orders, rather than
parameter orders, appear in \eqref{eq:finite-edge-v95}.
\end{example}
